%!TEX root = TheoretischePhysikA.tex

\section{Vorlesung: Zweikörperstoß und Lineare Algebra}

\begin{beispiel}[Zweikörperstoß]
Man betrachte zwei Körper, die zu einem Zeitpunkt zusammenstoßen: Die initialen und finalen Zustände werden wie folgt beschrieben:
\figbox
Es gilt 
\begin{equation}
  \begin{split}
    &\text{Impulserhaltung:}\quad \mbf{p}_1+\mbf{p}_2 = \mbf{p}_1'+\mbf{p}_2'\quad \text{und} \\
    &\text{Energieerhaltung:}\quad \frac{1}{2m_1}\mbf{p}_1\cdot \mbf{p}_1 +\frac{1}{2m_2}\mbf{p}_2\cdot \mbf{p}_2 = \frac{1}{2m_1}\mbf{p}_1'\cdot \mbf{p}_1' +\frac{1}{2m_2}\mbf{p}_2'\cdot \mbf{p}_2' + Q,
  \end{split}
\end{equation}
wobei $Q$ die deformations-Energie darstellt. Im Fall, dass $Q=0$ wird die Kollision als \emph{elastisch} bezeichnet, sonst als \emph{inelastisch}. Wir unterteilen inelastische Kollisionen in zwei Klassen, nämlich: endotherm $Q>0$ und exotherm $Q<0$.

Die Erhaltungsgleichungen liefern $4$ Gleichungen für $13$ Variablen ($12$ Impulskomponenten und $Q$). Aus gegebenen Anfangsimpulsen $\mbf{p}_1$ und $\mbf{p}_2$ können die Impulse nach dem Stoß und $Q$ (insgesamt $7$ Variablen) nicht eindeutig aus Energie- und Impulserhaltung bestimmt werden. Dafür müsste das Wechselwirkungspotential bestimmt sein, oder weitere Information über die Positionen der Teilchen vorliegen.
\end{beispiel}

\begin{beispiel}[Spezialfall: elastischer Stoß in $1D$]
Hierbei wird $Q=0$. Die Erhaltungsgleichungen lauten
\begin{equation}
	\begin{split}
		p_1+p_2 &= p_1'+p_2'\\
		\frac{p_1^2}{2m_1} + \frac{p_2^2}{2m_2} &= \frac{{p_1'}^2}{2m_1} + \frac{{p_2'}^2}{2m_2}.
	\end{split}
\end{equation}
Dieses Gleichungssystem bildet zwei Gleichungen für $2$ Unbekannte $p_1'$ und $p_2'$. Formen wir diese Gleichungen um, erhalten wir
\begin{equation}
	\begin{split}
		m_1(p_1-p_1') &= m_2(p_2'-p_2) \\
		m_1(p_1^2-{p_1'}^2) &= m_2({p_2'}^2-p_2^2).
	\end{split}
\end{equation}
Es gibt eine triviale Lösung, nämlich $p_1 = p_1'$ und $p_2 = p_2'$, die keiner Kollision entspricht. Gesucht wird allerdings eine nicht-triviale Lösung. Wir erhalten eine solche Lösung indem wir die Gleichungen teilen:
\begin{equation}
	\begin{split}
		\frac{p_1^2-{p_1'}^2}{p_1-p_1'} &= \frac{{p_2'}^2-p_2^2}{p_2'-p_2} \\
		p_1+p_1' &= p_2'+p_2.
	\end{split}
\end{equation}
Also
\begin{equation}
	p_2' = p_1+p_1'-p_2.
\end{equation}
Nach einsetzen in die Impulserhaltungsgleichung erhalten wir
\begin{equation}
	m_1p_1 + m_2p_2 = m_1p_1' + m_2(p_1+p_1'-p_2).
\end{equation}
Also
\begin{equation}
	\begin{split}
		p_1' &= \frac{m_1-m_2}{m_1+m_2}p_1 + \frac{2m_1}{m_1+m_2}p_2 \\
		p_2' &= \frac{2m_1}{m_1+m_2}p_1 + \frac{m_2-m_1}{m_1+m_2}p_2.
	\end{split}
\end{equation}
Insbesondere, für $m_1=m_2$ gilt $p_1'=p_2$ und $p_2'=p_1$: Die Stoßpartner tauschen ihre Impulse aus.

Es gibt ein alternatives Vorgehen: Man betrachte das System in einem Bezugssystem $\Sigma$, das sich in gleichförmiger Bewegung mit dem Schwerpunkt mitbewegt. 
Wir definieren dazu die \emph{Schwerpunktsposition}, \emph{Schwerpunktsgeschwindigkeit} und \emph{Schwerpunktsimpuls}:
\begin{equation}
	\begin{split}
		R &= \frac{m_1}{m_1+m_2}r_1 + \frac{m_2}{m_1+m_2}r_2 \\
		v &= \frac{m_1}{m_1+m_2}v_1 + \frac{m_2}{m_1+m_2}v_2 \\
		P &= (m_1+m_2)v = p_1+p_2. 
	\end{split}
\end{equation}
Es gilt
\begin{equation}
v = \frac{1}{m_1+m_2}p_1 + \frac{1}{m_1+m_2}p_2.
\end{equation}
\end{beispiel}

\subsection{Lineare Algebra: Matrizenrechnung}

Ein \emph{Vektor} $\mbf{v}\in \mathbb{R}^n$ ist eine Spalte von $n$ reellen Zahlen
\begin{equation}
	\mbf{v} = 
	\begin{pmatrix}
		v_1 \\ v_2 \\ \vdots \\ v_n
	\end{pmatrix}.
\end{equation}
Es gelten die folgenden Rechenoperationen
\begin{enumerate}
	\item Multiplikation mit Skalar $a\in\mathbb{R}$
	\begin{equation}
		a\mbf{v} = 
	\begin{pmatrix}
		av_1 \\ av_2 \\ \vdots \\ av_n
	\end{pmatrix}.
	\end{equation}
	\item Addition
	\begin{equation}
		\mbf{v}+\mbf{w} = \begin{pmatrix}
		v_1 \\ v_2 \\ \vdots \\ v_n
	\end{pmatrix} + \begin{pmatrix}
		w_1 \\ w_2 \\ \vdots \\ w_n
	\end{pmatrix} =
	\begin{pmatrix}
		v_1+w_1 \\ v_2+w_2 \\ \vdots \\ v_n+w_n
	\end{pmatrix}.
	\end{equation}
	\item Skalarprodukt
	\begin{equation}
		\mbf{v}\cdot\mbf{w} = \begin{pmatrix}
		v_1 \\ v_2 \\ \vdots \\ v_n
	\end{pmatrix} \cdot \begin{pmatrix}
		w_1 \\ w_2 \\ \vdots \\ w_n
	\end{pmatrix} =
	\sum_{j=1}^n v_jw_j.
	\end{equation}
\end{enumerate}

Wir führen ein verallgemeinertes Konzept des Vektoren ein: Eine \emph{Matrix} $\mbf{A}$ mit $m$ Reihen und $n$ Spalten definiert man als die folgende Anordnung von $n\cdot m$ reellen Zahlen:
\begin{equation}
	\mbf{A} \equiv \begin{pmatrix}
		A_{11} & A_{12} & A_{13} & \cdots & A_{1n} \\
		A_{21} & A_{22} & A_{23} & \cdots & A_{2n} \\
		\vdots & \vdots & &  \ddots  & \vdots \\
		A_{m1} & A_{m2} & A_{m3} & \cdots & A_{mn} \\
	\end{pmatrix}.
\end{equation} 
Zum Beispiel, ein $2\times 3$ Matrix besteht aus $2\cdot 3$ Zahlen angeordnet in $3$ Spalten und $2$ Reihen:
\begin{equation}
	\mbf{A} = \begin{pmatrix}
		1 & 3 & 2 \\ 1 & 9 & 0
	\end{pmatrix} = \begin{pmatrix}
		A_{11} & A_{12} & A_{13} \\ A_{21} & A_{22} & A_{23}
	\end{pmatrix}. 
\end{equation}
Das $j$-$k$ Element von $\mbf{A}$ wird $[\mbf{A}]_{jk}$ oder $A_{jk}$ geschrieben und steht in der $k$-ten Spalte und der $j$-ten Zeile. Zum Beispiel $[\mbf{A}]_{22} = 9$.

Es gelten die folgenden Rechenoperationen
\begin{enumerate}
	\item Multiplikation mit Skalar
	\begin{equation}
	a\mbf{A} \equiv \begin{pmatrix}
		aA_{11} & aA_{12} & aA_{13} & \cdots & aA_{1n} \\
		aA_{21} & aA_{22} & aA_{23} & \cdots & aA_{2n} \\
		\vdots & \vdots & &  \ddots  & \vdots \\
		aA_{m1} & aA_{m2} & aA_{m3} & \cdots & aA_{mn} \\
	\end{pmatrix}.
	\end{equation} 
	Also $[a\mbf{A}]_{jk} \equiv aA_{jk}$.
	\item Addition
	\item Matrixprodukt zwischen $n\times m$ Matrix $\mbf{A}$ und $m\times l$ Matrix $\mbf{B}$: das Matrixprodukt ergibt eine $n\times l$ Matrix $\mbf{C}$ mit Elementen
	\begin{equation}
		[\mbf{C}]_{jk} \equiv \sum_{a = 1}^m [\mbf{A}]_{ja}[\mbf{B}]_{ak}.
	\end{equation}
	Man kann das $j-k$-Element der Matrix $\mbf{C} = \mbf{A}\mbf{B}$ als das Skalarprodukt des $j$-ten Reihenvektors von $\mbf{A}$ und dem $k$-ten Spaltenvektor von $\mbf{B}$ betrachten.
\end{enumerate}

Das Matrixprodukt ist assoziativ, das heisst
\begin{equation}
	\mbf{A}\mbf{B}\mbf{C} = (\mbf{A}\mbf{B})\mbf{C} = \mbf{A}(\mbf{B}\mbf{C}).
\end{equation}
Das Matrixprodukt ist aber nicht kommutativ im Allgemeinen: es ist möglich, dass
\begin{equation}
	\mbf{A}\mbf{B} \not= \mbf{B}\mbf{A}.
\end{equation}

\begin{beispiel}
	Seien 
	\begin{equation}
		\mbf{A} = \begin{pmatrix}
			1 & 2 \\ 3 & 4
		\end{pmatrix} \quad\text{und}\quad \mbf{B} = \begin{pmatrix}
			5 & 6 & 7 \\ 8 & 9 & 10
		\end{pmatrix}
	\end{equation}
	zwei Matrizen. Wie lautet das Matrixprodukt $\mbf{A}\mbf{B}$?
\end{beispiel}
\begin{lösung}
	\begin{equation}
		\begin{split}
			\mbf{A}\mbf{B} &= \begin{pmatrix}
			1 & 2 \\ 3 & 4
		\end{pmatrix}\begin{pmatrix}
			5 & 6 & 7 \\ 8 & 9 & 10
		\end{pmatrix} \\
		&= \begin{pmatrix}
			1\times 5 + 2\times 8 & 1\times 6 + 2\times 9 & 1\times 7 + 2\times 10  \\
			3\times 5 + 4\times 8 & 3\times 6 + 4\times 9 & 3\times 7 + 4\times 10 
		\end{pmatrix} \\
		&= \begin{pmatrix}
			21 & 24 & 27 \\
			47 & 54 & 61
		\end{pmatrix} = \mbf{C}.
		\end{split}
	\end{equation}
\end{lösung}

\subsection{Lösen von linearen inhomogenen Gleichungssystemen}
Man betrachte ein inhomogenes System aus $2$ linearen Gleichungen mit $2$ Unbekannten $x_1$ und $x_2$:
\begin{equation}\label{eq:lin2}
	\begin{split}
		ax_1 + d x_2 &= y_1 \\
		cx_1 + b x_2 &= y_2
	\end{split}.
\end{equation}
Hierbei sind $a,b,c,d$ sowie $y_1,y_2$ gegeben. Wir können dieses Gleichungssystem als Matrixgleichung beschreiben
\begin{equation}
	\begin{pmatrix}
		a & d \\ c & b
	\end{pmatrix} \begin{pmatrix}
		x_1\\ x_2
	\end{pmatrix} = \begin{pmatrix}
		y_1\\ y_2
	\end{pmatrix}.
\end{equation} 
Diese Gleichung, welche wir als 
\begin{equation}
	\mbf{A}\mbf{x} = \mbf{y}
\end{equation}
schreiben, enthält dieselbe Information wie das System (\ref{eq:lin2}). Um diese Matrixgleichung zu lösen, bilden wir die \emph{inverse Matrix}
\begin{equation}
	\mbf{A}^{-1} \equiv \frac{1}{ab-cd} \begin{pmatrix}
	b & -d \\ -c & a
	\end{pmatrix}.
\end{equation}
Diese Matrix erfüllt die Relation
\begin{equation}
	\mbf{A}^{-1}\mbf{A} = \frac{1}{ab-cd} \begin{pmatrix}
	b & -d \\ -c & a
	\end{pmatrix} \begin{pmatrix}
		a & d \\ c & b
	\end{pmatrix} = \frac{1}{ab-cd} \begin{pmatrix}
	ba-dc & 0 \\ 0 & -cd+ab
	\end{pmatrix}  = \begin{pmatrix}
	1 & 0\\ 0 & 1
	\end{pmatrix} 
\end{equation}
Wir definieren dazu die sogenannte \emph{Einheitsmatrix}:
\begin{equation}
	\mathbb{I}_2\equiv\begin{pmatrix}
	1 & 0\\ 0 & 1
	\end{pmatrix}.
\end{equation}

Multiplizieren wir die Matrixgleichung mit $\mbf{A}^{-1}$ links mit dem Matrixprodukt, erhalten wir
\begin{equation}
	\begin{split}
		\mbf{A}^{-1}\mbf{A}\mbf{x} &= \mbf{A}^{-1}\mbf{y} \\
		\mathbb{I}_2 \mbf{x} &= \mbf{A}^{-1}\mbf{y} \\ 
		\begin{pmatrix}
	1 & 0\\ 0 & 1
	\end{pmatrix} \begin{pmatrix}
		x_1\\ x_2
	\end{pmatrix} &= \frac{1}{ab-cd} \begin{pmatrix}
	b & -d \\ -c & a
	\end{pmatrix} \begin{pmatrix}
		y_1 \\ y_2
	\end{pmatrix} \\
	\begin{pmatrix}
		x_1\\ x_2
	\end{pmatrix} &= \frac{1}{ab-cd} \begin{pmatrix}
		by_1 - dy_2 \\ -cy_1 + ay_2
	\end{pmatrix}.
	\end{split}
\end{equation}
Zusammengefasst: die Inverse Matrix $\mbf{A}^{-1}$ liefert die Lösung eines linearen Gleichungssystems.

Die Inverse Matrix 
\begin{equation}\label{eq:22inverse}
	\mbf{A}^{-1} \equiv \frac{1}{ab-cd} \begin{pmatrix}
	b & -d \\ -c & a
	\end{pmatrix}
\end{equation}
existiert nur dann, wenn $ab-cd\not=0$. Diese Größe wird als die \emph{Determinante der Matrix}
\begin{equation}
	\det(\mbf{A}) \equiv \det \begin{pmatrix}
		a & d \\ c & b
	\end{pmatrix} = ab-cd
\end{equation}
bezeichnet.


Dieses Verfahren wird folgendermaßen verallgemeinert. Ein inhomogenes Gleichungssystem aus $n$ Gleichungen mit $n$ Unbekannten 
\begin{equation}
	\begin{split}
		A_{11} x_1 + A_{12} x_2 + \cdots + A_{1n}x_n &= y_1 \\
		A_{21} x_1 + A_{22} x_2 + \cdots + A_{2n}x_n &= y_2 \\
		&\vdots  \\
		A_{n1} x_1 + A_{n2} x_2 + \cdots + A_{nn}x_n &= y_n 
	\end{split}
\end{equation}
kann wie folgt als Matrixgleichung beschrieben werden:
\begin{equation}\label{eq:linsysmat}
	\mbf{A}\mbf{x} = \mbf{y},
\end{equation}
wobei
\begin{equation}
	\mbf{A} = \begin{pmatrix}
		A_{11} & A_{12} & A_{13} & \cdots & A_{1n} \\
		A_{21} & A_{22} & A_{23} & \cdots & A_{2n} \\
		\vdots & \vdots & &  \ddots  & \vdots \\
		A_{n1} & A_{n2} & A_{n3} & \cdots & A_{nn} \\
	\end{pmatrix}, \quad \mbf{x} = \begin{pmatrix}
		x_1 \\ x_2 \\ \vdots \\ x_n
	\end{pmatrix}
\end{equation}
und
\begin{equation}
	\mbf{y} = \begin{pmatrix}
		y_1 \\ y_2 \\ \vdots \\ y_n
	\end{pmatrix}.
\end{equation}

Angenommen, es existiert eine Matrix $\mbf{A}^{-1}$ die die relation
\begin{equation}
	\mbf{A}^{-1} \mbf{A} = \begin{pmatrix}
		1 & 0 & 0 & \cdots & 0 \\
		0 & 1 & 0 & \cdots & 0 \\
		\vdots & \vdots & &  \ddots  & \vdots \\
		0 & 0 & 0 & \cdots & 1
	\end{pmatrix} \equiv \mathbb{I}_n
\end{equation}
erfüllt. Dann wird man eine Lösung des Gleichungssystems erhalten, indem man (\ref{eq:linsysmat}) von links mit $\mbf{A}^{-1}$ multipliziert:
\begin{equation}
 	\mbf{A}^{-1}\mbf{A}\mbf{x} = \mathbb{I}_n \mbf{x} = \mbf{x} = \mbf{A}^{-1}\mbf{y}.
\end{equation} 

Es stellt sich heraus, dass eine solche inverse Matrix $\mbf{A}^{-1}$ nur dann existiert, wenn $\det(\mbf{A})\not=0$. Hierbei definiert man die Determinante einer Matrix durch
\begin{equation}\label{}
	\det(\mbf{A}) = \sum_{j=1}^n (-1)^{j+k} [\mbf{A}]_{jk} \det(\mbf{A}_{jk}),
\end{equation}
wobei $\mbf{A}_{jk}$ die $(n-1)\times (n-1)$-Untermatrix von $\mbf{A}$ ist, die durch Streichen der $j$-ten Zeile und $k$-ten Spalte entsteht. Die Unterdeterminanten 
\begin{equation}
	\det (\mbf{A}_{jk})
\end{equation}
werden auch als \emph{Minoren} von $\mbf{A}$ bezeichnet. Die Zahlen
\begin{equation}
	\tilde {a}_{jk}=(-1)^{j+k}\cdot \det (\mbf{A}_{jk})
\end{equation}
heißen auch \emph{Kofaktoren} von $\mbf{A}$ und bilden als Matrix zusammengefasst die \emph{Kofaktormatrix} 
\begin{equation}
[\text{cof}(\mbf{A})]_{jk}={\tilde {a}}_{jk}.
\end{equation}
Die \emph{Transponierte} der Kofaktormatrix wird auch \emph{Adjunkte} $\text{adj}(\mbf{A})$ von $\mbf{A}$ genannt. Mit der Adjunkten hat die Inverse einer Matrix dann die explizite Darstellung
\begin{equation}\label{eq:inverse}
\mbf{A}^{-1}={\frac {1}{\det(\mbf{A})}}\cdot \text{adj}(\mbf{A}).
\end{equation}

In diesem Kurs werden nur $2\times 2$ und $3\times 3$ Matrizen untersucht. Wir haben schon die inverse Matrix im $2\times 2$ Fall aufgeschrieben. Der $3\times 3$ Fall ist etwas umständlicher: wenn man die Gleichung (\ref{eq:inverse}) für eine beliebige $3\times 3$ Matrix 
\begin{equation}
	\mbf{A} = \begin{pmatrix}
		a&b&c \\
		d&e&f \\
		g&h&i \\
	\end{pmatrix}
\end{equation}
aufschreibt, erhält man
\begin{equation}
	\mbf{A}^{-1} = \frac{1}{\det(\mbf{A})}\begin{pmatrix}
		A & D & G \\
		B & E & H \\
		C & F & I \\		
	\end{pmatrix},
\end{equation}
wobei
\begin{equation}
	\begin{split}
		A &= (ei-fh)\\
		B &= -(di-fg)\\
		C &= (dh-eg)\\
		D &= -(bi-ch)\\
		E &= (ai-cg)\\
		F &= -(ah-bg)\\
		G &= (bf-ce)\\
		H &= -(af-cd)\\
		I &= (ae-bd)\\
	\end{split}
\end{equation}
und
\begin{equation}
	\det(\mbf{A}) = aA+bB+cC.
\end{equation}